% Options for packages loaded elsewhere
\PassOptionsToPackage{unicode}{hyperref}
\PassOptionsToPackage{hyphens}{url}
\documentclass[
]{article}
\usepackage{xcolor}
\usepackage[margin=1in]{geometry}
\usepackage{amsmath,amssymb}
\setcounter{secnumdepth}{-\maxdimen} % remove section numbering
\usepackage{iftex}
\ifPDFTeX
  \usepackage[T1]{fontenc}
  \usepackage[utf8]{inputenc}
  \usepackage{textcomp} % provide euro and other symbols
\else % if luatex or xetex
  \usepackage{unicode-math} % this also loads fontspec
  \defaultfontfeatures{Scale=MatchLowercase}
  \defaultfontfeatures[\rmfamily]{Ligatures=TeX,Scale=1}
\fi
\usepackage{lmodern}
\ifPDFTeX\else
  % xetex/luatex font selection
\fi
% Use upquote if available, for straight quotes in verbatim environments
\IfFileExists{upquote.sty}{\usepackage{upquote}}{}
\IfFileExists{microtype.sty}{% use microtype if available
  \usepackage[]{microtype}
  \UseMicrotypeSet[protrusion]{basicmath} % disable protrusion for tt fonts
}{}
\makeatletter
\@ifundefined{KOMAClassName}{% if non-KOMA class
  \IfFileExists{parskip.sty}{%
    \usepackage{parskip}
  }{% else
    \setlength{\parindent}{0pt}
    \setlength{\parskip}{6pt plus 2pt minus 1pt}}
}{% if KOMA class
  \KOMAoptions{parskip=half}}
\makeatother
\usepackage{graphicx}
\makeatletter
\newsavebox\pandoc@box
\newcommand*\pandocbounded[1]{% scales image to fit in text height/width
  \sbox\pandoc@box{#1}%
  \Gscale@div\@tempa{\textheight}{\dimexpr\ht\pandoc@box+\dp\pandoc@box\relax}%
  \Gscale@div\@tempb{\linewidth}{\wd\pandoc@box}%
  \ifdim\@tempb\p@<\@tempa\p@\let\@tempa\@tempb\fi% select the smaller of both
  \ifdim\@tempa\p@<\p@\scalebox{\@tempa}{\usebox\pandoc@box}%
  \else\usebox{\pandoc@box}%
  \fi%
}
% Set default figure placement to htbp
\def\fps@figure{htbp}
\makeatother
\setlength{\emergencystretch}{3em} % prevent overfull lines
\providecommand{\tightlist}{%
  \setlength{\itemsep}{0pt}\setlength{\parskip}{0pt}}
\usepackage{bookmark}
\IfFileExists{xurl.sty}{\usepackage{xurl}}{} % add URL line breaks if available
\urlstyle{same}
\hypersetup{
  hidelinks,
  pdfcreator={LaTeX via pandoc}}

\author{}
\date{\vspace{-2.5em}}

\begin{document}

\section{逐笔成交量化作业口播稿}\label{ux9010ux7b14ux6210ux4ea4ux91cfux5316ux4f5cux4e1aux53e3ux64adux7a3f}

\begin{quote}
建议时长：7-8 分钟。括号内是提示，不需要逐字念。
\end{quote}

\subsection{0:00-0:40
开场与总体进度}\label{ux5f00ux573aux4e0eux603bux4f53ux8fdbux5ea6}

大家好，下面我汇报一下这次逐笔成交量化作业的完成情况。这次作业的主线是：先把
300 只股票、302
个交易日的逐笔成交数据降采样为日频和分钟频，然后构建和评价因子，再做 IC
加权的 Top 10 组合回测，最后用 LSTM
预测下一分钟涨跌。目前主任务和两个拓展都已完成，并且全量验证结果是
PASS，4 项单元测试也全部通过。

\subsection{0:40-2:00 数据处理}\label{ux6570ux636eux5904ux7406}

第一部分是数据降采样。原始字段包括成交时间、成交价格、成交量和主买主卖标志。我输出了
11
个基础字段，包括开高低收、成交量、成交笔数、成交额，以及主买主卖的量和金额。

在时间索引上，程序保留了开盘集合竞价、上下午连续竞价，以及 15
点收盘分钟，每天一共 253
行。对于无成交分钟，价格用前一分钟收盘填充；日内开始的空白用前一天收盘填充；量和金额则填零。这样可以避免用当日未来价格向前填充。

另外，题目补丁要求用复权因子修正价格。我对开高低收使用``原始价除以
100，再乘复权因子''，但成交额仍按实际成交价计算。最终日频每个字段都是
302 乘 300 的宽表，分钟频一共 3322 张表。

\subsection{2:00-3:20
因子构建与评价}\label{ux56e0ux5b50ux6784ux5efaux4e0eux8bc4ux4ef7}

第二部分是因子。我共构建了 4
个日频因子。第一个是题目给的示例因子，也就是 1、5、10、20
日成交额均值之间的标准差，再做对数变换。另外三个自建因子是 5
日主动买卖不平衡、10 日日内振幅和 5 日反转。

在评价前，我每天对因子横截面做 1\% 和 99\% 缩尾，再做 z-score
标准化。预测目标是下一交易日收盘到收盘收益。评价指标包括
IC、IR、ICIR、Rank IC、Rank IR、Rank ICIR，以及五分组和 Q5 减 Q1
多空收益。

从结果看，成交额多窗口因子的 Rank IC 是负
0.1070，绝对值最大，说明它是稳定的反向因子。5 日反转因子的 Rank IC 是正
0.0558，Q5 减 Q1 累计收益是
50.53\%。主动买卖不平衡因子接近于零，单独解释力比较弱。

\subsection{3:20-5:10
组合回测和前视偏差}\label{ux7ec4ux5408ux56deux6d4bux548cux524dux89c6ux504fux5dee}

第三部分是回测。按题目设定，初始资金为 1000 万，每天按综合信号选前 10
只股票等权持有，卖出手续费为万分之五，买入费用不计。我同时计算了收盘到收盘，和次日开盘到下一个开盘两种持有口径。

这里最重要的是前视偏差。题目的字面表述是用当日因子和次日收益计算当日
IC，然后立即用来加权。但在当日调仓时，次日收益还没发生，所以这个方案使用了未来信息。我保留它作为错误对照，但不把它当成策略结果。对照结果出现了
2200\% 到 2480\% 的总收益，这种异常高收益正是数据泄漏的警报。

修正方案是每个交易日只使用此前 20 日 IC
的均值，而且整体滞后一天。在这个可信版本下，收盘到收盘的总收益是
65.75\%，夏普率 2.291，最大回撤负 10.55\%。次日开盘到开盘的总收益是
23.72\%，夏普率 0.998，最大回撤负
23.00\%。后一个口径更保守，也说明策略对实际成交时点很敏感。

\subsection{5:10-6:35 LSTM 拓展}\label{lstm-ux62d3ux5c55}

第四部分是 LSTM。当前实验使用股票 000012 的前 60 个交易日，用过去 20
分钟预测下一分钟是否上涨。输入特征是对数收益、分钟振幅、对数成交量、主动买卖不平衡和日内时间位置。模型是单层
PyTorch LSTM，隐藏维度 16，再接一个线性二分类头。

为了避免时序泄漏，数据按时间顺序切分为 70\% 训练、15\% 验证和 15\%
测试，标准化参数只用训练集估计。训练使用
Adam、二元交叉熵和梯度裁剪，保留验证损失最低的检查点。

测试集有 2097 个样本，准确率 58.99\%，平衡准确率 51.05\%，AUC 是
0.5654。不过，多数类基线准确率是
59.56\%，比模型还略高。所以这一部分的结论是：端到端流程已经跑通，但当前信号很弱，不能说已经有可交易的预测能力。

\subsection{6:35-7:30
验证、局限和下一步}\label{ux9a8cux8bc1ux5c40ux9650ux548cux4e0bux4e00ux6b65}

最后说一下验证和后续工作。全量审计检查了日频和分钟频表的维度、对齐、有限数、高开低收关系和空值填充，也检查了
90600 行 K 线长表、因子和回测结果行数、每次是否持有 10 只股票，以及
PyTorch 模型检查点。当前结果全部通过。

下一步主要有三点。第一，回测中还要加入滑点、冲击成本、涨跌停和停牌限制，并进一步处理收盘信号同价成交的乐观假设。第二，因子可以增加行业和市值中性化。第三，LSTM
需要扩展到多股票、更长时间，并用 walk-forward 方式做严格样本外验证。

\subsection{7:30-7:50 结束}\label{ux7ed3ux675f}

总结来说，这次作业已经完成从原始逐笔数据到因子、回测和深度学习的完整链路。最重要的经验是，结果不能只看收益高低，还必须检查信号在当时是否真的可知。这也是我把前视偏差对照和可交易修正版分开展示的原因。我的汇报就到这里，谢谢。

\end{document}
